\documentclass[11pt]{article}
\usepackage{setspace}
\usepackage{fullpage}
\usepackage{amsmath}
\usepackage{multirow,array}
\usepackage[usestackEOL]{stackengine}
\usepackage[protrusion=true,expansion=true]{microtype} % Better typography
\usepackage{graphicx} % Required for including pictures
\usepackage{mdframed}
\usepackage{wrapfig} % Allows in-line images
\usepackage{csquotes}
\usepackage{booktabs}
\usepackage{morefloats}
\usepackage[toc,page]{appendix}
\usepackage{caption}
\usepackage{subcaption}
\usepackage{lipsum}
\usepackage{enumitem}
\usepackage{threeparttable}
\usepackage{pgf, tikz}
\usetikzlibrary{arrows, automata}
\usepackage[linguistics]{forest}
\usepackage{pgfplots}
\usepackage{etex}
\usepackage{natbib}
\setcitestyle{aysep={}}
\usepackage{mathpazo} % Use the Palatino font
\usepackage[T1]{fontenc} % Required for accented characters
\linespread{1.1} % Change line spacing here, Palatino benefits from a slight increase by default
%\usepackage{times}
\usepackage{url}
\usepackage{color}
\usepackage{pdfpages}
\usepackage{fancyhdr}
\usepackage{multicol}

\usepackage{rotating}
\usepackage{pdflscape}
\usepackage{xcolor}
\usepackage{import}
%\newcommand*{\Data}{../Data}% 
\usepackage{verbatim}
\usepackage{hyperref}
\usepackage{comment}
\usepackage[margin=1in]{geometry}
\usepackage{amsmath,amssymb}
\usepackage{setspace}
\setlength{\parskip}{6pt}
\setlength{\parindent}{0pt}

\title{Task 1: \\ The 4T Will Not Be Televised, It Will Be Livestreamed: Media, Censorship, and Narrative Control in Mexico"}
\author{Eduardo Zago}
\date{\today}

\begin{document}
\maketitle

\section{Quick overview}

This project builds on ideas I have been developing since my first semester. I initially worked on the empirical strategy in a proposal for Social Media and Politics, and later, in Game Theory II, I began designing a formal model that could help clarify the mechanisms at play and potentially serve as the basis for a structural estimation. I realize this may give the impression that I am presenting work that is already completed, but that is not the case. I simply want to be transparent that parts of the idea were previously developed in coursework. So I will now set up and explain the things that are done, and what I plan to do for this class.

I believe this class provides an ideal setting to advance the project in a more integrated and rigorous way. First, I would like to refine and complete the theoretical framework, developing a model that clearly addresses the research question—or a feasible version of it. Second, I aim to connect this model more explicitly to the empirical strategy, building on the data I have already collected and further strengthening that component. Third, I hope to improve the empirical design and, ideally, obtain preliminary results by the end of the term. Finally, if feasible, I would like to explore the possibility of structural estimation, using the model to generate broader insights.

\subsection{Motivation and theory}

How do governments maintain narrative control in democracies without sacrificing credibility? I wish to study the daily press conference of the Mexican President (\textit{La Mañanera}) as a dynamic Bayesian persuasion mechanism. I argue that the government faces a fundamental trade-off: minimizing dissent protects the immediate narrative but can render the communication uninformative, inducing rational disengagement by audiences, a result previously identified by \citep{Censor-Pratt-AER, GehlbachSonin2014}. 

Given the structure of the media market in Mexico--where the federal government has been a major purchaser of advertising since the 1990s \citep{Times, FourthState-Lawson-Book}--media outlets face a clear strategic tradeoff in their interactions with the executive. Attending presidential conferences and posing aligned or non-confrontational questions may help secure favorable treatment from the government, including greater access to public advertising funds. At the same time, such alignment may attract private advertising from audiences that support the incumbent administration.

However, as argued by \cite{GehlbachSonin2014}, repeated compliance carries long-term costs. When outlets consistently refrain from questioning the executive, their coverage provides little new information. Over time, audiences may disengage because the content becomes predictable and uninformative, ultimately reducing viewership and revenue.

Conversely, oppositional media outlets face a different set of incentives. A more critical stance may jeopardize access to public advertising funds, particularly if government advertising is allocated strategically. Yet critical coverage may attract private advertising by appealing to audiences seeking scrutiny and accountability, or by benefiting from the attention generated by controversial reporting. As a result, outlets must balance short-term financial incentives linked to government resources against longer-term reputational and market considerations.

\subsection{Research question and empirical strategy}

I argue that to restore credibility, the administration introduced a randomized seat lottery in 2022, effectively committing to a signal structure that allows for stochastic dissent.

This random assignment of seating--specifically, which outlets are placed in the first row, a position that (as I will show empirically) increases the likelihood of being selected to ask a question--provides a source of exogenous variation that allows me to identify two causal effects:

\begin{itemize}
\item What is the reduced-form short-term effect of media outlets asking critical questions to the government on their revenue?
\item Does the government respond strategically to critical outlets--for example, by reducing public advertising spending directed toward them?
\item What is the long-run equilibrium dissent allowed by the government and optimal for the media outlets?
\end{itemize}

\subsubsection*{Design 1 (ideal data): Instrumental Variables with observed attendance and seating}
\textbf{Data in the ideal scenario.} For each day $t$ and outlet $i$: (a) attendance indicator; (b) seating assignment, including a dummy $Z_{it}=\mathbf{1}\{\text{front row}\}$; (c) who asked; (d) transcript text to classify whether a question is \emph{critical}; (e) outcomes

\textbf{Variables.} Let $D_{it}$ be the number (or share) of \emph{critical questions} asked by outlet $i$ on day $t$ (constructed by supervised text methods on the question portion of the transcript); let $\Delta Y_{i,t+h}$ be an outcome

\textbf{Equations.} With outlet and (week?) fixed effects,
\[
\underbrace{D_{it}}_{\text{critical Q's}}\;=\;\pi\,Z_{it}\;+\;\lambda\,\text{Attend}_{it}\;+\;X'_{it}\gamma\;+\;\alpha_i\;+\;\tau_t\;+\;\nu_{it},
\]
\[
\underbrace{\Delta Y_{i,t+h}}_{\text{ad spending response}}\;=\;\beta\,\widehat{D}_{it}\;+\;X'_{it}\theta\;+\;\alpha_i\;+\;\tau_t\;+\;\varepsilon_{it}.
\]
\textbf{Interpretation.} $\beta$ is the LATE of an extra \emph{critical question} on ad spending for \emph{compliers to seating}. 

\textbf{Identification.} (1) \emph{Relevance}: front-row increases recognition $\Rightarrow$ more questions and higher probability a question is critical. (2) \emph{Independence}: conditional on attending and day fixed effects (to absorb ``hot'' news shocks), seating is as-if random among attendees. (3) \emph{Exclusion}: seating affects ad spending only through $D_{it}$. (4) \emph{Monotonicity}: no outlet becomes \emph{less} likely to ask when seated up front.

%------------------------------------------------------------

\subsubsection*{(NEW) Design 2 (no seating assignment): Difference in Difference with intermittent treatment}
\textbf{Data in the this scenario.} For each day $t$ and outlet $i$: (a) attendance indicator; (c) who asked; (d) transcript text to classify whether a question is \emph{critical}; (e) outcome, (f) media alignment

\textbf{Variables.} Let $D_{it}$ be a dummy that is 1 if outlet $i$ asked a question at time $t$; let $\Delta Y_{i,t+h}$ be the change in advertising allocated to $i$ between $t$ and $t{+}h$.

\textbf{Equations.} With outlet and time (weeks?) fixed effects,
\[
\underbrace{\Delta Y_{i,t+h}}_{\text{ad spending response}}\;=\;\beta\,D_{it}\;+\;\lambda\,\text{Attend}_{it}\;+\;X'_{it}\theta\;+\;\alpha_i\;+\;\tau_t\;+\;\varepsilon_{it}.
\]

\textbf{Interpretation.} $\beta$ is the ATT of a \emph{question} on ad spending conditional on attendance. Then, to understand the consequences of aligned vs. oppositional media, I could look at heterogeneous treatment effects by running the same specification for each set of aligned or oppositional media.

\textbf{Identification.} Using \cite{LiuEtAl} estimator and framework (1) \emph{Strict exogeneity}: past outcomes do not directly affect current treatment assignment (satisfied if truly there is randomization) (2) \emph{no carry-over effects}: treatment affects the contemporaneous outcome only, not future outcomes. Can be tested, and given the recurrence of the interaction, randomization and dynamic of media articles this should also hold. 


\subsection*{Current feasibility and data status}

At present I only have the official transcripts for all mañaneras. From these I can (i) \emph{observe who asked} each day and (ii) often infer \emph{where the asker sat} from the video/transcript structure (not always, and the government has not been helpful on getting this data, which is why the second design might be better). I have the list of which outlet attended each day $Attend_{it}$. 

For outcomes, at the moment, I have the total government spending on advertising by outlet, the day the contract for each ad was signed and the amount of each contract. This semester I plan to obtain the following outcomes: 

\begin{itemize}
    \item Number of views, comments, likes on each outlet's Youtube videos and Twitter profiles (daily).
    \item The number of times each outlet was searched on Google (Google Trends data).
    \item Will probably need more time but all the articles posted by each newspaper.
    \item If we are able, get actual private advertising data from Nielsen or other company (has been done on \citep{MediaCaptures-Zeidl-Econometrica})
\end{itemize}

\subsection{Model and structural (?) estimation}

I would say the main issue of this empirical strategy is that it allows to only causally identify the immediate reduced-form effect on dissent/alignment. For the second design is pretty clear this is the case as the no carry over assumption restricts me from analyzing temporal effects. However, the theory that I just presented suggests that this strategic interaction is dynamic, and therefore a model could be used to first pinned down the exact mechanisms at play and hopefully, through structural estimation or other method, being able to speak about the dynamic effects. For Game Theory 2 I developed a simple model where the government decides on the level of allowed dissent:

\textbf{Setup: The Government (Sender) and the Public (Receiver)}

\begin{itemize}
    \setlength\itemsep{0.8em}
    \item \textbf{State of the World:} $\theta \in \{G, B\}$ (Good, Bad).
    \item \textbf{Prior Belief:} $q = \Pr(\theta = G)$.
    \item \textbf{Public's Action:} $a \in \{0, 1\}$ (Disengage, Support).
    \begin{itemize}
        \item The public supports ($a=1$) iff belief $\Pr(G \mid m) \ge \tau$ (credibility threshold).
    \end{itemize}
    \item \textbf{Assumption (Skepticism):} $q < \tau$. 
    \begin{itemize}
        \item \footnotesize{The government starts with a credibility deficit. Without new information, the public disengages.}
    \end{itemize}
    \item \textbf{Signal (The Broadcast):} $m \in \{S, O\}$.
    \begin{itemize}
        \item $S$: \textbf{Scripted} (Only favorable questions).
        \item $O$: \textbf{Open} (Critical dissent is visible).
    \end{itemize}
\end{itemize}


Why Total Censorship Fails

Suppose the government chooses \textbf{Total Control}: always produce a Scripted broadcast ($S$), regardless of the state.

\begin{itemize}
    \item \textbf{Information Structure:}
    \[ \Pr(S \mid G) = 1, \quad \Pr(S \mid B) = 1 \]
    
    \item \textbf{Public's Inference (Bayes' Rule):}
    Since the message is uninformative, the posterior equals the prior:
    \[ \Pr(G \mid S) = q \]
    
    \item \textbf{Outcome:}
    Since $q < \tau$ (Skepticism), the public chooses $a=0$ (Disengage).
    
    \item \textbf{Result:} 
    Total censorship yields zero political support. The broadcast is discounted as "pure propaganda." in line with \citep{Censor-Pratt-AER} conclusions
\end{itemize}


The Optimal Solution: Interior Dissent

The government maximizes the probability of Support ($a=1$) by choosing a signal structure $\beta = \Pr(S \mid B)$ (how often to hide the Bad state).

\textbf{Optimal Strategy:}
\begin{enumerate}
    \item Always show Scripted if things are Good: $\Pr(S \mid G) = 1$.
    \item Choose $\beta < 1$ such that the "Scripted" signal is \textit{just credible enough}:
    \[ \Pr(G \mid S) = \frac{q \cdot 1}{q \cdot 1 + (1-q) \cdot \beta} = \tau \]
\end{enumerate}

\textbf{The Interior Solution:}
\[ \beta^* = \frac{q}{1-q} \frac{1-\tau}{\tau} < 1 \]

 To maximize support, the government \textbf{must allow some dissent} ($\Pr(O \mid B) > 0$). The optimal level of censorship is interior, not absolute.


The idea for this class would be to iterate on this model in the following way: 

\begin{itemize}
    \item Add media outlets as strategic agents. Decide on whether asking a critical or aligned question or the severity of their question.
    \item Model this as a repeated interaction game (few periods). First, government starts, then media starts, etc.
    \item Model this as a dynamic Bayesian persuasion mechanism (Markov model)
    \item Learn about structural estimation and maybe estimate the model, or set it up in a way that maps to the data and I am able to estimate it.
\end{itemize}



%------------------------------------------------------------


\clearpage
\small
\bibliographystyle{chicago-apsr-ph}
\typeout{}

\singlespacing

\bibliographystyle{chicago-apsr-ph}
\bibliography{literature}


\end{document}